\chapter{Marco teórico}\label{cap:marco}

En este capítulo se describen los conceptos y términos que se utilizarán en el desarrollo de esta tesis.

\section{Fundamentos de recuperación de información}\label{sec:marco-fundamentos}

\subsection{El problema de la recuperación de información}\label{subsec:marco-problema}
% - Definición informal: consulta + corpus → ranking de documentos
% - Definiciones formales: consulta, documento, corpus, relevancia, ranking como permutación
% - El sistema de RI como función de puntuación y el orden que induce
% - Relevancia binaria vs. graduada (anticipo para §2.5)

\subsection{El principio de ordenación por probabilidad (PRP)}\label{subsec:marco-prp}
% - Enunciado: ordenar por P(relevante | d, q) es óptimo bajo sus supuestos
% - Consecuencia: los modelos de RI como estimadores (puente hacia el BIM)

\section{Modelos de recuperación de primera etapa}\label{sec:marco-primera-etapa}

\subsection{Modelos léxicos probabilísticos: del BIM a BM25}\label{subsec:marco-bm25}
% - Bolsa de palabras y TF-IDF como heurística clásica
% - BIM y peso RSJ: derivación desde el PRP
% - 2-Poisson: saturación de la frecuencia de término
% - Normalización por longitud del documento
% - BM25 canónico: fórmula final, papel de k₁ y b

\subsection{Representaciones dispersas aprendidas (SPLADE)}\label{subsec:marco-splade}
% - Misma idea que BM25 (vector sobre el vocabulario) con pesos aprendidos
% - SPLADE: proyección sobre vocabulario vía MLM, expansión aprendida, regularización de dispersión
% - Lectura conceptual: «BM25 aprendido» → léxico moderno de esta tesis

\subsection{Representaciones densas y codificadores duales}\label{subsec:marco-densos}
% - Hipótesis distribucional; de embeddings estáticos a contextuales
% - Tokenización y autoatención: lo mínimo del Transformer (encoder bidireccional)
% - Codificador dual: definición formal (dos torres + similitud), pooling
% - Qué captura lo semántico frente a lo léxico (motiva la complementariedad)
% - Mención breve: variantes guiadas por instrucciones (mE5-instruct) y entrenamiento contrastivo (InfoNCE)

\section{Reordenamiento neuronal}\label{sec:marco-reranking}

\subsection{La arquitectura en cascada: recuperación y reordenamiento}\label{subsec:marco-cascada}
% - Por qué en etapas: el costo prohíbe aplicar modelos pesados a todo el corpus
% - Primera etapa (recall) vs. segunda etapa (precisión); el recall inicial como cota

\subsection{Cross-encoders}\label{subsec:marco-crossencoders}
% - Definición: consulta y documento concatenados, puntuación conjunta
% - Interacción completa → más precisión; costo por par → no escala
% - Taxonomía pointwise / pairwise / listwise (listwise enlaza al cap. 3)

\section{Fusión de rangos}\label{sec:marco-fusion}

\subsection{Formalización}\label{subsec:marco-fusion-formal}
% - Ranking como permutación; notación formal
% - Equivalencia de rango bajo transformaciones monótonas → justifica la fusión de rangos
% - Fusión de puntuaciones vs. de rangos: el problema de las escalas incompatibles
% - Analogía con teoría de elección social (sistemas = votantes, documentos = candidatos)

\subsection{Métodos basados en puntuaciones (CombMNZ)}\label{subsec:marco-combmnz}
% - Normalización de puntuaciones como prerrequisito
% - CombSUM y CombMNZ; el factor MNZ que premia la presencia en varias listas

\subsection{Métodos posicionales (BordaFuse, RRF, ISR)}\label{subsec:marco-posicionales}
% - BordaFuse: puntos por posición, conexión con el conteo de Borda
% - RRF: 1/(k+r), papel del parámetro k
% - ISR (Inverse Square Rank): definición (peso 1/r²) y propiedades

\subsection{Comparaciones por pares y modelos de usuario (Condorcet, RBC)}\label{subsec:marco-pares}
% - Condorcet: mayoría por pares, grafo de preferencias y manejo de ciclos
% - RBC (Rank-Biased Centroids): modelo de usuario con profundidad de escaneo; fusión como utilidad esperada

\section{Evaluación en recuperación de información}\label{sec:marco-evaluacion}

\subsection{Métricas de evaluación}\label{subsec:marco-metricas}
% - P@k y Recall@k: precisión y cobertura a profundidad fija
% - MAP y MRR
% - nDCG: ganancia graduada con descuento; por qué es la métrica primaria

\subsection{Pruebas de significancia estadística}\label{subsec:marco-significancia}
% - El problema: promedios de métrica vs. variabilidad por consulta; comparaciones pareadas por consulta
% - Los cinco candidatos: t pareada, Wilcoxon, signo, bootstrap y permutación (randomization)
% - Evidencia empírica consolidada: t pareada y permutación mantienen el error tipo I nominal; bootstrap sesgado hacia p-valores pequeños; Wilcoxon y signo poco fiables para diferencias de medias (anticipa la elección del cap. 4)
% - Errores tipo I, II y III; pruebas de una y dos colas
% - Corrección por comparaciones múltiples (Bonferroni / Tukey HSD)
